% 神经网络损失景观：2D等高线 + 多鞍点/局部极小/平台 + 平坦 vs 尖锐极小对比框
\documentclass[tikz,border=4pt]{standalone}
\usepackage{ctex}
\usepackage{amsmath}
\usetikzlibrary{calc, arrows.meta, decorations.pathmorphing, backgrounds, fit}

\begin{document}
\begin{tikzpicture}[scale=1.0, thick]

% ── 左侧：2D 损失等高线示意 ──────────────────────────
\begin{scope}[xshift=0cm]
  % 背景：大范围等高线（用椭圆模拟）
  \draw[gray!30, thin] (0,0) ellipse (3.5cm and 2.5cm);
  \draw[gray!40, thin] (0,0) ellipse (2.8cm and 2.0cm);
  \draw[gray!50, thin] (0,0) ellipse (2.1cm and 1.5cm);
  \draw[gray!60, thin] (0,0) ellipse (1.4cm and 1.0cm);

  % 高损失区域中的多个鞍点（用 ×标记）
  \node[red, font=\large] at (-2.5, 1.2) {$\times$};
  \node[below right, font=\footnotesize, red] at (-2.5, 1.2) {鞍点};

  \node[red, font=\large] at (2.2, 1.5) {$\times$};
  \node[below, font=\footnotesize, red] at (2.2, 1.5) {鞍点};

  \node[red, font=\large] at (1.5, -1.8) {$\times$};
  \node[above right, font=\footnotesize, red] at (1.5, -1.8) {鞍点};

  % 局部极小（低损失区域）
  \filldraw[blue] (-0.8, 0.3) circle (2pt);
  \node[below left, font=\footnotesize, blue] at (-0.8, 0.3) {局部极小};

  % 全局极小（最低点）
  \filldraw[black!70] (0.5, -0.4) circle (3pt);
  \node[below right, font=\footnotesize, black!70] at (0.5, -0.4) {全局极小};

  % 平台区域（梯度接近零）
  \draw[gray!70, dashed, thick] (-1.5, -0.8) ellipse (0.6cm and 0.3cm);
  \node[below, font=\footnotesize, gray!70] at (-1.5, -0.8) {平台区};

  % 标题
  \node[font=\small\bfseries] at (0, -3.1) {损失景观（2D 等高线示意）};
  \node[font=\footnotesize] at (0, -3.5) {参数空间 $\theta_1$};
  \node[font=\footnotesize, rotate=90] at (-4.2, 0) {参数空间 $\theta_2$};

  % 轴
  \draw[->, gray] (-3.8, 0) -- (3.8, 0) node[right, font=\footnotesize] {$\theta_1$};
  \draw[->, gray] (0, -2.8) -- (0, 2.8) node[above, font=\footnotesize] {$\theta_2$};
\end{scope}

% ── 右侧：平坦 vs 尖锐极小对比 ──────────────────────
\begin{scope}[xshift=8.5cm]

  % 坐标轴
  \draw[->, gray] (-0.3, 0) -- (6.2, 0) node[right, font=\footnotesize] {$\theta$};
  \draw[->, gray] (0, -0.3) -- (0, 3.2) node[above, font=\footnotesize] {$\mathcal{L}$};

  % ── 平坦极小（宽碗形）──
  \draw[blue, very thick]
    plot[smooth, tension=0.8]
    coordinates {(0.2, 3.0) (0.8, 1.8) (1.5, 0.6) (2.5, 0.1) (3.5, 0.6) (4.2, 1.8) (4.8, 3.0)};
  \filldraw[blue] (2.5, 0.1) circle (2.5pt);
  \node[below, font=\footnotesize, blue] at (2.5, -0.1) {平坦极小 $\theta^*_\text{flat}$};

  % 平坦极小标注框
  \draw[blue, dashed, thin] (1.0, 0.0) -- (1.0, 0.55);
  \draw[blue, dashed, thin] (4.0, 0.0) -- (4.0, 0.55);
  \draw[blue, <->, thin] (1.0, 0.45) -- (4.0, 0.45) node[midway, above, font=\footnotesize, blue] {宽（泛化好）};

  % 训练/测试损失接近的标注
  \node[font=\footnotesize, blue, align=center] at (2.5, 2.5) {训练损失 $\approx$\\测试损失};

  % ── 尖锐极小（窄碗形）──
  \begin{scope}[xshift=0cm, yshift=0cm]
    % 尖锐极小在同一坐标系，用红色
    % 用分段曲线表示窄尖的形状
    \draw[red, very thick]
      plot[smooth, tension=1.2]
      coordinates {(0.4, 3.0) (0.8, 2.2) (1.2, 1.0) (1.5, 0.2) (1.8, 1.0) (2.2, 2.2) (2.6, 3.0)};
    \filldraw[red] (1.5, 0.2) circle (2.5pt);
    \node[below, font=\footnotesize, red] at (1.5, -0.1) {尖锐极小 $\theta^*_\text{sharp}$};

    % 尖锐极小宽度
    \draw[red, dashed, thin] (1.1, 0.0) -- (1.1, 0.35);
    \draw[red, dashed, thin] (1.9, 0.0) -- (1.9, 0.35);
    \draw[red, <->, thin] (1.1, 0.25) -- (1.9, 0.25) node[midway, above, font=\footnotesize, red] {窄};

    % 测试损失远高于训练
    \node[font=\footnotesize, red, align=center] at (1.5, 2.5) {测试损失\\$\gg$ 训练损失};
  \end{scope}

  % 图例
  \draw[blue, very thick] (3.5, 2.2) -- (4.3, 2.2);
  \node[right, font=\footnotesize, blue] at (4.3, 2.2) {平坦极小};
  \draw[red, very thick] (3.5, 1.8) -- (4.3, 1.8);
  \node[right, font=\footnotesize, red] at (4.3, 1.8) {尖锐极小};

  \node[font=\small\bfseries] at (3.0, -0.8) {平坦 vs 尖锐极小};
\end{scope}

\end{tikzpicture}
\end{document}
